\section{Vanishing twist and the complete static base}\label{rig4:sta:section}

We continue with the equality data and the fields supplied by
Theorem~\ref{rig4:adj:causal-field}. The stationary metric
\[
 \mathbf g=-u^2dt^2+g_{ij}(dx^i+X^i dt)(dx^j+X^j dt)
 \quad\text{on }\R\times\operatorname{int}\Omega
\]
is smooth because $u>0$ in the interior. Its Killing field is
$\xi=\partial_t$. Recall that
\[
 N=u^2-|X|_g^2\ge0,\qquad
 \Ric_{\mathbf g}(\xi,\cdot)=0,
 \qquad \Ric_{\mathbf g}=0\quad\text{where }N>0.
\]
Moreover, $N\to1$ at the end, and $N>0$ on a full interior collar
of $S$. Thus the interior zero set
\[
 Z=\{x\in\operatorname{int}\Omega:N(x)=0\}
\]
is compactly contained in $\operatorname{int}\Omega$. No regularity of
$Z$ as a subset is assumed. In particular, it may have interior or
infinitely many components.

\subsection{The Killing current near a null set}

On the open set $\mathcal P=\{N>0\}\subset\operatorname{int}\Omega$, put
\begin{equation}\label{rig4:sta:base-variables}
 \begin{aligned}
 h_b&=g+N^{-1}X^\flat\otimes X^\flat,
 & C_b&=g^{-1}-u^{-2}X\otimes X=h_b^{-1},\\
 A_b&=N^{-1}X^\flat,
 & F&=dA_b,\qquad \lambda=\sqrt N .
 \end{aligned}
\end{equation}
Here $X^\flat$ is formed with $g$. Direct substitution gives
\begin{equation}\label{rig4:sta:stationary-decomposition}
 \mathbf g=-N(dt-A_b)^2+h_b.
\end{equation}
The tensor $C_b$ extends smoothly and nonnegatively through $Z$, although
$h_b$ and $A_b$ need not extend there. For a covector $\alpha$ write
$|\alpha|_{C_b}^2=C_b^{ij}\alpha_i\alpha_j$.

\begin{lemma}[Killing-current identities]\label{rig4:sta:current-identities}
Let $\mathcal G_{\alpha\beta}=\nabla^{\mathbf g}_\alpha\xi_\beta$.
On the interior of $\Omega$ one has
\begin{equation}\label{rig4:sta:norm-equation}
 u^{-1}\operatorname{div}_g(uC_b\,dN)
       =-2|\mathcal G|_{\mathbf g}^2.
\end{equation}
On $\mathcal P$ the antisymmetric tensor
\begin{equation}\label{rig4:sta:killing-current}
 Q^{ij}=uN C_b^{ik}C_b^{j\ell}F_{k\ell}
 \qquad\text{satisfies}\qquad \nabla^g_iQ^{ij}=0.
\end{equation}
The norm in \eqref{rig4:sta:norm-equation} is the Lorentzian contraction of
both tensor indices; it need not be nonnegative.
\end{lemma}

\begin{proof}
The Killing identities and $\Ric_{\mathbf g}(\xi,\cdot)=0$ give
\[
 \nabla^{\mathbf g,\alpha}\mathcal G_{\alpha\beta}=0,
 \qquad
 \square_{\mathbf g}\bigl(\mathbf g(\xi,\xi)\bigr)
       =2|\mathcal G|_{\mathbf g}^2.
\]
For a stationary scalar function the wave operator is
$u^{-1}\operatorname{div}_g(uC_b\,d\,\cdot\,)$: the spacetime volume
density is $u\sqrt{\det g}$ and its spatial inverse-metric block is
$C_b$. Since $\mathbf g(\xi,\xi)=-N$, this proves
\eqref{rig4:sta:norm-equation}, including at $Z$ where all its entries are
smooth.

For the second identity set $\theta=dt-A_b$. On $\mathcal P$,
\[
 \xi^\flat=-N\theta,
 \qquad
 \mathcal G=\tfrac12d\xi^\flat
            =\tfrac12(-dN\wedge\theta+NF).
\]
The inverse-metric dual of $\theta$ is proportional to $\partial_t$,
so the term containing $\theta$ has zero spatial-spatial component after
both indices are raised. Consequently
\[
 \mathcal G^{ij}=\tfrac N2 C_b^{ik}C_b^{j\ell}F_{k\ell}.
\]
The stationary divergence identity, with density $u\sqrt{\det g}$,
now gives \eqref{rig4:sta:killing-current}. In converting coordinate
divergence to $g$-covariant divergence, the additional Christoffel term
vanishes by antisymmetry.
\end{proof}

\begin{lemma}[A transverse logarithmic estimate]\label{rig4:sta:log-estimate}
Suppose $Z\ne\varnothing$. On a neighborhood of $Z$ with compact closure
in the interior, let
\[
 B=|X|_g,\qquad e=X/B,\qquad D=e^\perp,
\]
and denote by $d_DN$ the $g$-orthogonal projection of $dN$ to $D$.
This neighborhood can be chosen so that $u$ and $B$ are bounded below
by positive constants. For every smooth cutoff $\zeta$ compactly
supported there,
\begin{equation}\label{rig4:sta:null-log-energy}
 \sup_{\varepsilon>0}
 \int u\zeta^2\frac{|dN|_{C_b}^2}{(N+\varepsilon)^2}\,dV_g
 <\infty.
\end{equation}
In particular, $d_D\log N$ is square-integrable on the positive set
in every smaller neighborhood of $Z$.
\end{lemma}

\begin{proof}
At a point of $Z$, $B=u>0$. Compactness provides the positive lower
bounds after shrinking a neighborhood. On that neighborhood the
eigenvalues of $C_b$ relative to $g$ are $1$ on $D$ and $N/u^2$ in the
$e$ direction. Thus
\begin{equation}\label{rig4:sta:degenerate-norm}
 |dN|_{C_b}^2=|d_DN|_g^2+\frac N{u^2}(eN)^2.
\end{equation}
Smoothness and nonnegativity of $N$ also give
\begin{equation}\label{rig4:sta:nonnegative-gradient}
 |dN|_g^2\le C N.
\end{equation}
For completeness, work in a slightly larger compact interior
neighborhood, where the Hessian is bounded. Taylor's inequality along
a short geodesic in the direction $-\nabla N$ gives
$0\le N-t|dN|+Ct^2/2$. Increasing $C$ so that $t=|dN|/C$ stays
in this neighborhood proves \eqref{rig4:sta:nonnegative-gradient}.

Let $n=u^{-1}(\xi-X)$ be the future unit normal to the original
slice. In a local orthonormal frame $n,e,e_a$, with the $e_a$ in $D$,
differentiating $\mathbf g(\xi,\xi)=-N$ gives
\[
 \mathcal G_{in}=-\frac{N_i/2+B\mathcal G_{ie}}u
 \quad\text{for each spatial index }i.
\]
As
$|\mathcal G|_{\mathbf g}^2
 =2\sum_{i<j}\mathcal G_{ij}^2-2\sum_i\mathcal G_{in}^2$,
we obtain the exact expansion
\begin{equation}\label{rig4:sta:killing-norm-bound}
 \begin{split}
 -2|\mathcal G|_{\mathbf g}^2
   ={}&\frac{|dN|_g^2}{u^2}
       +\frac{4B}{u^2}\sum_aN_a\mathcal G_{ae}
       -\frac{4N}{u^2}\sum_a\mathcal G_{ae}^2
       -4\sum_{a<b}\mathcal G_{ab}^2\\
   \le{}& C\bigl(N+|d_DN|_g\bigr).
 \end{split}
\end{equation}
All coefficients of $\mathcal G$ are bounded here. This estimate is
pointwise and does not require a global frame for $D$.

Write $f=-2|\mathcal G|_{\mathbf g}^2$ and test
$\operatorname{div}_g(uC_b\,dN)=uf$ against
$\zeta^2/(N+\varepsilon)$. There is no singular boundary in this test:
Equation~\eqref{rig4:sta:norm-equation} is smooth across $Z$, and
$\varepsilon>0$. Integration by parts gives
\begin{equation}\label{rig4:sta:log-test}
 \begin{split}
 \int u\zeta^2\frac{|dN|_{C_b}^2}{(N+\varepsilon)^2}\,dV_g
 ={}&\int\frac{u\zeta^2 f}{N+\varepsilon}\,dV_g\\
 &+2\int\frac{u\zeta\langle d\zeta,dN\rangle_{C_b}}
                  {N+\varepsilon}\,dV_g.
 \end{split}
\end{equation}
In the first term use \eqref{rig4:sta:killing-norm-bound},
$N/(N+\varepsilon)\le1$, and
$|d_DN|\le|dN|_{C_b}$. Young's inequality absorbs its derivative
term into, say, one quarter of the left side. The second integral is
handled by Cauchy's inequality for the nonnegative tensor $C_b$ and
another quarter of that side. The remaining terms are bounded by a
constant times
$\int u(\zeta^2+|d\zeta|_{C_b}^2)\,dV_g$, uniformly in
$\varepsilon$. This proves \eqref{rig4:sta:null-log-energy}. Fatou's lemma
and \eqref{rig4:sta:degenerate-norm} give the last assertion.
\end{proof}

\subsection{Removal of all twist boundary terms}

\begin{proposition}[Vanishing of the twist]\label{rig4:sta:twist-vanishes}
On all of $\mathcal P$, including every component, $F=dA_b=0$.
\end{proposition}

\begin{proof}
On the asymptotic end the constant limit of $A_b$ is the covector
$a_\infty=b^i\,dx^i$, because $N\to1$ and $g\to\delta$.
Choose a smooth function $f$ on $\Omega$ which equals $b^i x^i$
sufficiently far out and is zero off an end neighborhood disjoint
from $S$ and $Z$. Then
\[
 \beta=A_b-df,\qquad d\beta=F,
 \qquad \beta=O(r^{-q_0})\quad\text{at infinity}.
\]
For any nonnegative smooth function $\chi$ compactly supported in
$\mathcal P$, multiply \eqref{rig4:sta:killing-current} by $\chi\beta_j$
and integrate. Antisymmetry gives
\begin{equation}\label{rig4:sta:twist-test}
 \frac12\int_{\mathcal P}
   \chi uN C_b^{ik}C_b^{j\ell}F_{ij}F_{k\ell}\,dV_g
  =-\int_{\mathcal P}Q^{ij}(\partial_i\chi)\beta_j\,dV_g.
\end{equation}
The integrand on the left is nonnegative, and its quadratic form in
$F$ is positive definite at each point of $\mathcal P$.

The contraction in the right side has a cancellation that is needed
at the null loci. Lowering its remaining index with $g$, one finds
\begin{equation}\label{rig4:sta:current-contraction}
 \begin{split}
 g_{ik}Q^{kj}(A_b)_j
 &=\frac N u F_{ij}X^j\\
 &=\frac1u\left[
 X^j(dX^\flat)_{ij}
 -\frac{|X|^2N_i-X(N)X_i}{N}\right].
 \end{split}
\end{equation}
Indeed $C_b A_b=X/u^2$, and the projection in the other raised index
does not change $F_{ij}X^j$. This covector is orthogonal to $X$.
Where $X\ne0$, its last numerator is $B^2(d_DN)_i$.
The coordinate-free numerator in \eqref{rig4:sta:current-contraction}
will also be used where $X=0$.

We choose cutoffs near the three possible escape loci with disjoint
transition neighborhoods. If $Z=\varnothing$, the first cutoff below
is identically one.

\paragraph{The interior null set.}
Choose nested relatively compact interior neighborhoods of $Z$ on
which $u$ and $B$ are bounded below and
Lemma~\ref{rig4:sta:log-estimate} applies. A smooth cutoff $\eta$ is one on
the smaller neighborhood and supported in the larger. Let
$\vartheta:[0,\infty)\to[0,1]$ be smooth, zero on $[0,1]$ and one on
$[2,\infty)$, with bounded derivative, and set
\[
 \chi_{Z,\varepsilon}=1-\eta+\eta\vartheta(N/\varepsilon).
\]
For all sufficiently small $\varepsilon$, $N$ has a positive lower
bound on the support of $d\eta$, so differentiation of $\eta$ adds
no term. The only transition is $\varepsilon<N<2\varepsilon$, and
its derivative is bounded by $C|dN|/\varepsilon$.
Here $df=0$. Orthogonality in
\eqref{rig4:sta:current-contraction} replaces $dN$ by $d_DN$ in the
contraction. Since $N\asymp\varepsilon$ on the band and the smooth
term $X\mathbin{\lrcorner}dX^\flat$ is bounded, the absolute
contribution to the right side of \eqref{rig4:sta:twist-test} is at most
\begin{equation}\label{rig4:sta:null-cutoff-cost}
 C\int_{\{\varepsilon<N<2\varepsilon\}\cap\supp\eta}
       \bigl(1+|d_D\log N|_g^2\bigr)\,dV_g=o(1).
\end{equation}
The integrable majorant comes from
Lemma~\ref{rig4:sta:log-estimate}. The characteristic functions of these
bands tend pointwise to zero even at points of $Z$. Thus dominated
convergence proves the last assertion without an assumption on the
measure or structure of $Z$.

\paragraph{The asymptotic end.}
Let $\chi_{\infty,R}$ be one for $r\le R$, zero for $r\ge2R$,
with derivative bounded by $C/R$, and extend it by one over the
compact part. From the adjoint end estimates,
\[
 Q=O(r^{-1-q_0}),\qquad \beta=O(r^{-q_0}).
\]
The transition annulus has volume $O(R^4)$, so its contribution is
\begin{equation}\label{rig4:sta:end-cutoff-cost}
 O(R^{2-2q_0})=o(1),
\end{equation}
since $q_0>1$.

\paragraph{The original boundary.}
Let $s$ be original $g$-normal distance from $S$, increasing into
$\Omega$, and let $\chi_{S,\sigma}$ change from zero to one on
$\sigma<s<2\sigma$, with derivative bounded by $C/\sigma$.
Again $df=0$ near this collar. If $u|_S>0$, the boundary conclusions
of Theorem~\ref{rig4:adj:causal-field} give
\[
 dN=2\kappa X^\flat\quad\text{on }S,
 \qquad X=u\nu\quad\text{on }S,
 \qquad N\asymp s.
\]
In particular $d_DN=O(s)$, so
\eqref{rig4:sta:current-contraction} is bounded. Its contraction with the
normal $\nabla s$ is $O(s)$: it annihilates $X$, and
$\nabla s-X/u=O(s)$. Multiplication by $1/\sigma$ and integration
over a collar of volume $O(\sigma)$ therefore gives $O(\sigma)$.

If $u|_S=0$, the same theorem supplies smooth one-sided fields
$d,Y_2$ with
\begin{equation}\label{rig4:sta:zero-lapse-collar}
 u=sd,\qquad X=s^2Y_2,\qquad d|_S=\kappa,
 \qquad N=s^2(d^2-s^2|Y_2|^2).
\end{equation}
Thus $u\asymp s$, $N\asymp s^2$, $dX^\flat=O(s)$, and
$dN=O(s)$. Both terms inside the brackets of
\eqref{rig4:sta:current-contraction} are $O(s^3)$; its entire value is
$O(s^2)$. This gives a boundary cutoff cost $O(\sigma^2)$.

Now take
$\chi=\chi_{Z,\varepsilon}\chi_{S,\sigma}\chi_{\infty,R}$.
For positive parameters with sufficiently small $\varepsilon,\sigma$
and large $R$, it is a compactly supported smooth test on
$\mathcal P$. Multiplication by the other two factors does not
increase any of the three estimates. Therefore the absolute value
of the right side of \eqref{rig4:sta:twist-test} tends to zero as
$\varepsilon,\sigma\downarrow0$ and $R\to\infty$. Choose a sequence
of such parameters; its cutoffs tend pointwise to one on
$\mathcal P$. Fatou's lemma gives zero for the integral of the
nonnegative density on the left. Its continuity and pointwise
positive definiteness force $F=0$ everywhere on $\mathcal P$.
This integration is on an open manifold with compactly supported
tests and remains valid if $\mathcal P$ has infinitely many
components.
\end{proof}

\subsection{Static equations and the only finite-distance attachment}

Let $\mathcal U$ be the component of $\mathcal P$ containing the
distant asymptotic end. Since $F=0$ on all of $\mathcal P$, the
Poincar\'e lemma supplies a local primitive of $A_b$ on every
sufficiently small ball in any positive component. Replacing $t$ locally
by $T=t-a$ with $da=A_b$ turns
\eqref{rig4:sta:stationary-decomposition} into
$-\lambda^2dT^2+h_b$. Its Ricci components are
\[
 (\Ric_{\mathbf g})_{TT}=\lambda\Delta_{h_b}\lambda,
 \qquad
 (\Ric_{\mathbf g})_{ij}
      =(\Ric_{h_b})_{ij}-\lambda^{-1}(\Hess_{h_b}\lambda)_{ij}.
\]
These follow directly from the nonzero mixed Christoffel symbols
$\Gamma^T_{Ti}=\partial_i\log\lambda$ and
$\Gamma^i_{TT}=\lambda h_b^{ij}\partial_j\lambda$.
Vacuum on $\mathcal P$ therefore gives, on every positive component,
\begin{equation}\label{rig4:sta:static-equations}
 \lambda\Ric_{h_b}=\Hess_{h_b}\lambda,
 \qquad \Delta_{h_b}\lambda=0,
 \qquad \Scal_{h_b}=0.
\end{equation}
This conclusion uses local primitives only; no global time primitive
or topological conclusion has yet been used.

\begin{lemma}[Connectedness of the positive set]
\label{rig4:sta:connected-positive-set}
The set $\mathcal P$ is connected and equals $\mathcal U$. In particular,
the positive collar of $S$ belongs to $\mathcal U$.
\end{lemma}

\begin{proof}
The unique asymptotic end lies in $\mathcal U$ outside a compact set.
Any other component $V$ of $\mathcal P$ therefore has compact closure
in the original exterior $\Omega$, with its boundary included.
Its frontier lies in $S\cup Z$: an interior frontier point with $N>0$
would have a connected positive neighborhood meeting $V$ and hence
belong to $V$. The frontier is nonempty, since otherwise $V$ would
be both open and closed in the connected interior of $\Omega$, which
also contains $\mathcal U$.

The continuous function $\lambda=\sqrt N$ vanishes on this frontier
and is positive in $V$. It therefore attains a positive maximum at
an interior point of $V$. The strong maximum principle for
$\Delta_{h_b}\lambda=0$ makes $\lambda$ constant on $V$, contradicting
its zero frontier values. Thus there is no such component $V$.
The collar supplied by Theorem~\ref{rig4:adj:causal-field} lies in
$\mathcal P=\mathcal U$, as claimed.
\end{proof}

\begin{proposition}[Completeness and regular horizon attachment]
\label{rig4:sta:complete-base}
The component $\mathcal U$ satisfies $0<\lambda<1$. Every approach in
$\mathcal U$ to an interior zero of $N$ has infinite $h_b$-length.
Adjoining $S$ gives a complete manifold with smooth one-sided metric
$h_b$ and lapse $\lambda$, in the boundary structure specified below, and
\begin{equation}\label{rig4:sta:boundary-data}
 h_b|_{TS}=g|_{TS},\qquad \lambda|_S=0,
 \qquad \partial_{\nu_{h_b}}\lambda=\kappa,
 \qquad \mathrm{II}_{h_b}=0.
\end{equation}
Here $\nu_{h_b}$ points into the base exterior. The metric has an even
reflection and the lapse an odd reflection of class at least $C^2$.
More precisely:
\begin{enumerate}[label=(\roman*)]
\item If $u|_S>0$, then $N$ is a defining function in the original
smooth structure and $dN=2\kappa X^\flat$ on $S$. The base boundary
structure uses $\lambda=\sqrt N$ in place of $N$. In that structure
the reflected metric and signed lapse are smooth.
\item If $u|_S=0$, the one-sided base metric and lapse are smooth in
the original boundary structure. Their even and odd reflections in
base normal distance are $C^2$.
\end{enumerate}
In either case the base attachment is homeomorphic to the original
attachment of $S$. There are no other finite-distance boundary
attachments.
\end{proposition}

\begin{proof}
First consider an interior approach to $Z$. With the notation of
Lemma~\ref{rig4:sta:log-estimate}, Equation~\eqref{rig4:sta:current-contraction}
and $F=0$ give
\[
 B^2d_D\log N=X^j(dX^\flat)_{ij}\,dx^i,
 \qquad |d_D\log N|_g\le C.
\]
Together with \eqref{rig4:sta:nonnegative-gradient} this yields
\begin{equation}\label{rig4:sta:base-log-bound}
 |d\log N|_{h_b}^2
 =|d_D\log N|_g^2+\frac{(eN)^2}{u^2N}\le C
\end{equation}
on the positive set near $Z$. Along any curve approaching $Z$,
$\log N\to-\infty$, which is incompatible with finite length and
\eqref{rig4:sta:base-log-bound}.

The original space $(\Omega,g)$, with its boundary included, is a
complete locally compact length space and hence proper. Since
$h_b\ge g$ on $\mathcal U$, a finite-length base curve is $g$-Cauchy
and has a limit in $\Omega$. A limit in the interior with $N>0$ lies
in $\mathcal U$ and is a regular base point. A limit in $Z$ is
excluded by the preceding bound. Escape to the asymptotic end also
has infinite length. Thus the only possible finite-distance
attachment is $S$, whose positive collar belongs to $\mathcal U$
by Lemma~\ref{rig4:sta:connected-positive-set}.

Next, $\lambda$ tends to one at infinity, and it tends to zero at
any finite limiting boundary of $\mathcal U$ in the original
manifold. If $\lambda>1$ somewhere, continuity on the compact part
of the original space and its end limit give an attained interior
maximum. The strong maximum principle in
\eqref{rig4:sta:static-equations} would make $\lambda$ constant on
$\mathcal U$, contradicting its limit one. Hence $\lambda\le1$.
If equality held at an interior point, the same principle would
give $N\equiv1$ on $\mathcal U$. Such a component has no interior
frontier, because continuity would give $N=1>0$ at each frontier
point. It is therefore both open and closed in the connected
$\operatorname{int}\Omega$, so it is the entire interior. This
contradicts $N\to0$ at $S$. Thus $0<\lambda<1$.

It remains to construct and check the attachment of $S$.

\paragraph{Positive boundary lapse.}
Here $dN=2u\kappa\nu^\flat$ on $S$, with $u>0$, so $(N,y^A)$
are original local boundary coordinates. The identity
$dN=2\kappa X^\flat$ at $N=0$ and smooth division give
\[
 X^\flat=a\,dN+N\beta_A\,dy^A,
 \qquad a(0,y)=\frac1{2\kappa},
\]
where $a$ and $\beta_A$ are smooth in $(N,y)$. Passing to
$(\lambda,y)$, with $N=\lambda^2$, the base metric components are
\begin{equation}\label{rig4:sta:positive-lapse-base-collar}
 \begin{split}
 (h_b)_{\lambda\lambda}
       &=4\lambda^2g_{NN}+4a^2,\\
 (h_b)_{\lambda A}
       &=2\lambda(g_{NA}+a\beta_A),\\
 (h_b)_{AB}&=g_{AB}+\lambda^2\beta_A\beta_B.
 \end{split}
\end{equation}
All original coefficients on the right are evaluated at
$(\lambda^2,y)$. The first and third expressions are smooth even
functions of signed $\lambda$, and the middle one is smooth odd.
They therefore define a smooth metric invariant under
$(\lambda,y)\mapsto(-\lambda,y)$. At the boundary the cross terms
vanish and $(h_b)_{\lambda\lambda}=\kappa^{-2}$; the tangential
metric is $g|_{TS}$. Nondegeneracy holds on both sides after
restricting the collar. The reflection fixes $S$, so its second
fundamental form vanishes, and its inward unit normal is
$\kappa\partial_\lambda$ there. This proves
\eqref{rig4:sta:boundary-data}. The lapse is the signed coordinate
$\lambda$, hence is smooth and odd. Changes between these boundary
charts are smooth after the substitution $N=\lambda^2$, so the
construction is compatible on overlapping collars.

\paragraph{Zero boundary lapse.}
Equation~\eqref{rig4:sta:zero-lapse-collar} gives
\[
 h_b=g+
 \frac{s^2Y_2^\flat\otimes Y_2^\flat}{d^2-s^2|Y_2|^2},
 \qquad
 \lambda=s\sqrt{d^2-s^2|Y_2|^2}.
\]
Both are smooth one-sided in the original boundary structure,
$h_b=g$ at $S$, and $\partial_{\nu_{h_b}}\lambda=\kappa$.
Since $\Ric_{h_b}$ is smooth up to $S$, the first equation in
\eqref{rig4:sta:static-equations} extends continuously and gives
$\Hess_{h_b}\lambda=0$ there. Its tangential part equals
$\kappa\mathrm{II}_{h_b}$, proving that the second form vanishes.
Use base Gaussian coordinates $(\rho,y)$, so
$h_b=d\rho^2+h_{AB}(\rho,y)dy^A dy^B$. At $\rho=0$ we have
\[
 \partial_\rho h_{AB}=0,
 \qquad \lambda=0,
 \qquad \partial_\rho^2\lambda=0.
\]
Even reflection of $h_{AB}$ and odd reflection of $\lambda$ thus
match derivatives through order two. This proves the claimed $C^2$
reflection and all of \eqref{rig4:sta:boundary-data}.

These constructions give a positive one-sided smooth base metric
up to $S$. In the first case replacing $N\ge0$ by
$\lambda=\sqrt N\ge0$ is a homeomorphism; in the second case the
original structure is retained. The attachment therefore has the
same underlying topology as the original one. If a base Cauchy
sequence approaches $S$ in the original topology, the corresponding
regular base coordinates converge, so it converges in the attached
base metric as well. Combining this observation with the
finite-length argument above proves completeness with $S$ included.
\end{proof}
